\section{Maximum irredundant length and a normal module}
A generating family is irredundant if deleting any member leaves a
proper subgroup. Write $\mu(G)$ for the largest size of an irredundant
generating family of a finite group $G$, and $\mu'(G)$ for the largest
size of any independent family in it. These are maximum lengths,
not the minimum number of generators. For an abelian normal subgroup
$V\unlhd G$, let $\mu'_G(V)$ denote the largest size of a family
independent under generation by all $G$-conjugates.

\begin{lemma}[UB.NormalQuotient]\label{lem:quotient}
For finite $G$ and $N\unlhd G$,
$\mu(G)\leq\mu(G/N)+\mu'(N)$. If $N$ is abelian, the final term
can be replaced by $\mu'_G(N)$. In particular, if $V\unlhd G$ is
an elementary abelian minimal normal subgroup, then
\begin{equation}\label{eq:normalmodule}
 \mu(G)\leq\mu(G/V)+1.                                      \end{equation}
No splitting of the extension is needed.
\end{lemma}
\begin{proof}
Take an irredundant generating family $S$ of $G$, and an inclusion-minimal
subfamily $T$ whose images generate $G/N$. For $s\in S\setminus T$,
choose a word $w_s$ in $T$ with $sN=w_sN$, and put
$v_s=sw_s^{-1}\in N$. If one $v_s$ lies in the subgroup generated by
the others, substitution expresses $s$ using $S\setminus\{s\}$,
a contradiction. Since $|T|\leq\mu(G/N)$, the first bound follows.

When $N$ is abelian, $G=N\gen{T}$ implies that its conjugation action
on $N$ is already the action of $\gen{T}$. Thus an expression for
$v_s$ in conjugates of the other $v_t$ again substitutes to an
expression for $s$ without $s$. This proves the normal-module bound.
If $V$ is elementary abelian minimal normal, it is an irreducible
$\mathbb F_pG$-module. Any nonzero vector generates all of $V$ under
conjugation, so an independent family for this action has at most one
member. This proves \eqref{eq:normalmodule} for maximum irredundant length.
\end{proof}

This is Cameron--Cara's normal-quotient theorem [CC, Theorem 1.1]. Both the extremal and the high-rank applications use this one complete proof. We retain Whiston's bound that an independent family in $S_m$
has at most $m-1$ members, with equality only when it generates $S_m$.
In particular $\mu'(A_m)\leq m-2$. No classification of the C-groups
of $\mathrm{PSL}_2$ or $\mathrm{PGL}_2$ is used below.

\begin{corollary}\label{cor:affine}
If $G$ has a faithful primitive affine action with elementary abelian
translation subgroup $V$, then \eqref{eq:normalmodule} applies. In particular
$\mu(\mathrm{AGL}(4,2))\leq7$.
\end{corollary}
\begin{proof}
The translation module is irreducible: a proper invariant subspace
would give a nontrivial block system. Thus $V$ is minimal normal.
For $\mathrm{AGL}(4,2)$ the quotient is
$\mathrm{GL}(4,2)\cong A_8$, so Whiston gives $\mu(A_8)\leq6$.
The classical exceptional isomorphism is a retained identification;
the finite certificate also checks an explicit isomorphism for the
catalogue representation used below.
\end{proof}

\section{All transitive imprimitive groups of degree nine}
\begin{theorem}\label{thm:nine}
Every transitive subgroup of $S_3\wr S_3$ in its imprimitive action
on nine letters has $\mu(G)\leq5$. Consequently it cannot be a
maximal parabolic of a rank-seven C-group marking of $S_9$.
\end{theorem}
\begin{proof}
Let $K$ be the kernel of the action on the three blocks, and put
$A=K\cap A_3^3$. Then $A$ is an elementary abelian normal 3-subgroup,
viewed as an invariant subspace of $\mathbb F_3^3$. The quotient
$G/A$ has normal subgroup $W=K/A\leq C_2^3$ and transitive quotient
$B\leq S_3$, so $B=C_3$ or $S_3$.

Choose an element of $G$ inducing a three-cycle on the blocks.
On $A_3^3$ its conjugation action is a cyclic permutation of the
coordinates with coefficients $\varepsilon_i\in\{1,-1\}$.
The vector in any one coordinate is cyclic for this operator, whose
characteristic and minimal polynomials are
$X^3-\delta=(X-\delta)^3$, where $\delta=\varepsilon_1\varepsilon_2
\varepsilon_3\in\mathbb F_3$. Thus the full module is
$\mathbb F_3[X]/(X-\delta)^3$. Its invariant subspaces form a chain,
and each is cyclic. Any family of vectors in an invariant subspace
therefore has one vector whose cyclic span contains all the others.
Independence under the whole group implies independence under this
one element, so $\mu'_G(A)\leq1$.

The corresponding action on $C_2^3$ is ordinary coordinate permutation.
Under the three-cycle it is the direct sum of a trivial line and the
irreducible plane with polynomial $X^2+X+1$. Every invariant subspace
is a sum of a subset of these two nonisomorphic irreducibles. The two
polynomial projection operators show that a vector generates each
component on which it is nonzero. At most two vectors are therefore
needed from any given family to generate its cyclic submodule; a
family independent under the full group has at most two members.
Hence $\mu'_{G/A}(W)\leq2$.

Twice applying Lemma~\ref{lem:quotient} gives
\[
 \mu(G)\leq\mu(G/A)+1\leq\mu(B)+2+1\leq5.
\]
This applies to the maximum irredundant length, not just a selected
small generating set. A parabolic of rank six in a C-group has six
independent distinguished generators, contradicting the bound.
\end{proof}

\section{Independent vectors in a transitive binary module}
The following elementary observation is useful beyond the present
classification. It bounds irredundance of a particular family, rather
than merely providing some short generating family for the module.

\begin{lemma}\label{lem:simplequotients}
Let $M$ be a finite-dimensional module over a finite field, generated
irredundantly as a module by $v_1,\ldots,v_s$. There are simple quotients
$L_1,\ldots,L_s$ of $M$ and a surjection
$M\longrightarrow\bigoplus_{i=1}^s L_i$.
\end{lemma}
\begin{proof}
For each $i$, choose a maximal proper submodule $M_i$ containing all
$v_j$ with $j\ne i$. It does not contain $v_i$, since those vectors
generate $M$. Put $L_i=M/M_i$. Under the product quotient map, the
submodule generated by $v_i$ has zero coordinates except the $i$th,
where its image is the nonzero submodule $L_i$. Thus the image contains
every direct summand and is the whole direct sum.
\end{proof}

\begin{lemma}\label{lem:binary}
Let $P\leq S_m$ be transitive, where $m\geq4$, acting by coordinate
permutation on $V=\mathbb F_2^m$. For every invariant subspace
$W\leq V$, every $P$-independent family in $W$ has size at most $m-2$.
For $P=S_m$ the stronger bound two holds.
\end{lemma}
\begin{proof}
Let $M$ be the submodule generated by an independent family of size $s$.
If $M<V$ and $s\geq m-1$, Lemma~\ref{lem:simplequotients} forces
$\dim M=s=m-1$ and every simple quotient there has dimension one.
Over $\mathbb F_2$ a one-dimensional representation is trivial, so the
surjection is an isomorphism and $M$ is fixed pointwise by $P$.
But the fixed space of a transitive permutation group on $V$ is the
one-dimensional constant subspace, a contradiction.

If $M=V$, at most one of those simple quotients has dimension one:
otherwise the product surjection would give two linearly independent
maps from $V$ to the trivial module, whereas a $P$-invariant linear
functional has equal coefficients in all $m$ positions. Hence
$m\geq2s-1$, and $s\leq\lfloor(m+1)/2\rfloor\leq m-2$.

For the final assertion, if every vector in the family is constant,
one suffices. Otherwise choose a nonconstant vector $v$. Subtracting
its image under a transposition that exchanges a zero and a one gives
$e_i+e_j$. Conjugation then generates the entire even-weight subspace.
This subspace has codimension one, so at most one further member of the
family is needed. Independence therefore permits at most two.
\end{proof}

\begin{theorem}\label{thm:twoblocks}
Every transitive subgroup of $S_2\wr S_m$ in degree $2m$, $m\geq5$,
has maximum irredundant generating length at most $2m-4$.
\end{theorem}
\begin{proof}
Let $W\leq\mathbb F_2^m$ be the block kernel and $P\leq S_m$ the
transitive block image. The conjugation action on $W$ is coordinate
permutation. If $\mu(P)\leq m-2$, Lemmas~\ref{lem:quotient}
and~\ref{lem:binary} give $\mu(G)\leq2m-4$. Otherwise Whiston's
equality clause forces $P=S_m$. The sharper module bound gives
$\mu(G)\leq(m-1)+2=m+1\leq2m-4$.
\end{proof}

\begin{lemma}\label{lem:smallbase}
For $m\geq4$, every transitive subgroup of $S_3\wr S_m$ has
$\mu(G)\leq3m-4$. For $m\geq3$, every transitive subgroup of
$S_4\wr S_m$ has $\mu(G)\leq4m-4$.
\end{lemma}
\begin{proof}
For blocks of size three, write $K$ for the block kernel,
$A=K\cap A_3^m$, $W=K/A\leq\mathbb F_2^m$, and $P=G/K$.
If $P\ne S_m$, then $\mu(P)\leq m-2$, the binary-module cost
is at most $m-2$, and the ternary-module cost is at most its
ambient dimension $m$. This gives $3m-4$.

If $P=S_m$, the binary cost is at most two. The ternary cost is at
most $m-1$. Indeed, an independent family of $m$ vectors would span
all of $\mathbb F_3^m$, and Lemma~\ref{lem:simplequotients} would
identify this monomial module with a sum of one-dimensional modules.
Its action would then have abelian image. But its permutation image
on the coordinate axes is $S_m$, which is nonabelian. Thus the total
cost is at most $(m-1)+2+(m-1)=2m\leq3m-4$.
For three blocks of size three use Theorem~\ref{thm:nine} instead.

Now let the blocks have size four. The block-kernel projections are
conjugate. If they are proper subgroups of $S_4$, every independent
family in each projection has at most two members, by Whiston's
equality clause. Successive projections of the kernel and
Lemma~\ref{lem:quotient} bound its independent families by $2m$.
Consequently $\mu(G)\leq(m-1)+2m\leq4m-4$.

Otherwise every projection is $S_4$. Put $A=K\cap V_4^m$, where
$V_4$ is the normal Klein four subgroup of $S_4$. Then $G/A$ has a
faithful transitive imprimitive action with $m$ blocks of size three,
so the preceding bound gives $\mu(G/A)\leq3m-4$.
We claim that the normal-module cost of $A\leq\mathbb F_2^{2m}$
is at most $m$. The group $L=K/A$ is a subdirect subgroup of $S_3^m$.
Its normal subgroup $Q=L\cap C_3^m$ projects onto $C_3$ in every
coordinate: the commutator subgroup of $L$ already has this property.
Thus $Q$ fixes no nonzero vector in the full $2m$-dimensional module,
since the natural $C_3$ action on each two-dimensional coordinate
has no fixed vector. Averaging over the odd-order group $Q$ shows
that taking $Q$-fixed vectors is exact, so every subquotient of this
module also has zero $Q$-fixed space. In particular every simple
$G$-module quotient in Lemma~\ref{lem:simplequotients} has dimension
at least two. Hence a normal-independent family in $A$ has size
at most $\dim A/2\leq m$. The normal-module inequality now gives
$\mu(G)\leq(3m-4)+m=4m-4$.
\end{proof}

\section{A bounded primitive certificate}
For an integer $a>0$, let $\Omega(a)$ be its number of prime factors
with multiplicity. Every strictly ascending subgroup chain in a finite
group of order $a$ has at most $\Omega(a)$ strict inclusions, because
the product of its successive integral indices divides $a$. Thus the
C-rank is at most $\Omega(|G|)$.

The required program \texttt{primitive\_bound.g} enumerates the complete
GAP primitive-group catalogue in degrees 9--25, omitting $A_n,S_n$,
and reports exactly twelve groups for which this elementary bound
exceeds $n-4$. All remaining catalogue entries are thereby discharged.
The twelve exceptions are:
\begin{equation}\label{eq:primitiveexceptions}
\begin{array}{c|l}
9&3^2\!:\!QD_{16},\quad3^2\!:\!SL_2(3),\quad AGL_2(3),\quad
 PSL_2(8),\quad PSL_2(8)\!:\!3;\\
10&A_6\!:\!2,\quad S_6,\quad A_6\!.\!2,\quad A_6\!:\!2^2;\\
11&M_{11};\\12&M_{12};\\16&AGL_4(2).
\end{array}                                                   \end{equation}
The program verifies the precise subgroup and quotient identifications
used next. Its retained inputs are the completeness of this catalogue
and GAP's exact finite-group algorithms; it performs no search for
C-group markings and imports none of the source's reported C-ranks.

For the first two degree-nine groups the irreducible translation module
costs one and the quotient has at most four prime factors, giving five.
For $AGL_2(3)$ the quotient $GL_2(3)$ has Frattini subgroup of order two
and Frattini quotient $S_4$. Passing to a Frattini quotient preserves
irredundant generating families: a subset generates the whole group
exactly when its image generates the quotient. Thus $\mu(GL_2(3))
=\mu(S_4)\leq3$, and the affine group has $\mu\leq4$.
Every maximal subgroup of $PSL_2(8)$ has order 56, 18 or 14; these have
at most four prime factors, so a C-group marking has rank at most five.
The extension by 3 is not generated by its involutions: the program
checks that they generate only the index-three subgroup.

In degree ten the three groups other than $S_6$ have normal socle
$A_6$ and quotient of order two or four. Lemma~\ref{lem:quotient}
and $\mu'(A_6)\leq4$ give at most six generators; $S_6$ has at most
five. Corollary~\ref{cor:affine} treats degree sixteen.

For $M_{11}$ the maximal-subgroup orders are
$720,660,144,120,48$. The subgroup generated by all involutions in
each has at most six prime factors; in the order-720 case its order
is 360. Any C-group parabolic lying in such a maximal subgroup lies
in its involution-generated subgroup. Hence $M_{11}$ has C-rank at
most seven. For $M_{12}$ the maximal-subgroup orders are
\[
7920,7920,1440,1440,660,432,432,240,192,192,72.
\]
The first two are $M_{11}$. The order-1440 groups have normal $A_6$
with quotient order four, so their irredundant lengths are at most six.
Every remaining maximal subgroup has at most seven prime factors.
Deleting one distinguished generator and placing the resulting proper
parabolic in a maximal subgroup gives C-rank at most eight for $M_{12}$.
All these finite assertions are checked by the complete program.

For arbitrary degree $n\geq26$, retain Mar\'oti's Theorem 1.1 in its
original form, including the lower containment $A_m^\ell\leq G$ in
the product-action alternative. If $\ell\geq2$, the same abstract
group acts faithfully on $m\ell$ points, so Whiston bounds its
independent families by $m\ell-1\leq m^\ell-4\leq n-4$; the
inequality holds for $m\geq2$, $\ell\geq2$, $m^\ell\geq9$.
If $\ell=1$, the natural action case is excluded by $G\ne A_n,S_n$;
otherwise the $k$-subset action has $k\geq2$ and $m\leq n/2$,
so the same conclusion follows from $m-1\leq n-4$.
The Mathieu alternative has degree at most 24.

In the remaining alternative, Mar\'oti bounds the group order by
$M(n)=n\prod_{i=0}^{\lfloor\log_2n\rfloor-1}(n-2^i)$.
For $26\leq n\leq63$, exact integer arithmetic in the certificate
verifies $M(n)<2^{n-3}$. For $n\geq64$, put
$k=\lfloor\log_2n\rfloor\geq6$. Then
\[
 \log_2M(n)<(k+1)\log_2n<(k+1)^2\leq2^k-3\leq n-3.
\]
The middle inequality follows at $k=6$ and is preserved when $k$
increases by one. Since C-rank is an integer and is less than $n-3$,
it is at most $n-4$. We have proved the full primitive-parabolic
bound for all degrees, conditional only on the explicitly retained
Whiston and Mar\'oti theorems and the bounded exact catalogue inputs.

\section{A uniform simple-kernel bound}
\begin{lemma}\label{lem:scott}
Let $S$ be a finite nonabelian simple group. Every subdirect subgroup
of $S^m$ is a product of diagonal copies of $S$ on a partition of the
$m$ coordinates, with automorphisms allowed in the diagonal maps.
A strictly increasing chain of subdirect subgroups has at most $m-1$
strict steps.
\end{lemma}
\begin{proof}
Proceed by induction on $m$. Project onto the first $m-1$ coordinates.
The kernel in the last factor is normal in $S$, so is either $S$ or
trivial. In the first case the subgroup is the product of its projection
with $S$. In the second it is the graph of an onto homomorphism from
that projection to $S$. By induction the projection is a direct product
of simple diagonal factors. Their images are commuting normal subgroups
of $S$, so exactly one has nontrivial image, and its map is an
isomorphism. The last coordinate therefore joins that diagonal block.
This proves the description. In an inclusion of such subgroups the
larger group can only refine the partition. If the number of blocks
does not increase, their orders are equal and the groups coincide.
There are between one and $m$ blocks, proving the chain bound.
\end{proof}
This is the usual subdirect-product lemma, often called Scott's lemma;
its complete elementary proof is included here, rather than retained
as an unpriced extra premise.

\begin{lemma}\label{lem:simplekernel}
Suppose $N\unlhd G$ is a subdirect subgroup of $S^m$ for a finite
nonabelian simple group $S$, and conjugation by $G$ permutes the
coordinates transitively, using automorphisms on individual factors.
Then
\begin{equation}\label{eq:simplekernel}
 \mu(G)\leq\mu(G/N)+\mu'(S)+m-1.                    \end{equation}
\end{lemma}
\begin{proof}
Normalize an irredundant generating family as in
Lemma~\ref{lem:quotient}, obtaining $T\cup E$, with $T$ minimally
surjecting onto $G/N$ and $E\subset N$. The normalized extra elements stay independent relative to $\gen{T}$,
by the substitution argument, and no member of $T$ can be dropped since
its quotient image is indispensable. Thus this replacement preserves
irredundance. Write $T_0=\gen{T}$ and $N_0=T_0\cap N$.
The subgroup generated by all $T_0$-conjugates of $E$ is normalized by
$T_0$ and contains $E$, hence is normal in $G$. It follows that
$N$ is generated by $N_0$ and these conjugates.

Choose a minimal subset $B\subseteq E$ such that the projections of
$N_0$ and of all $T_0$-conjugates of $B$ generate $S$ in the first
coordinate. Then $|B|\leq\mu'(S)$. To see this, for each $b\in B$
choose an element in the projected subgroup generated by its conjugates
which does not belong to the subgroup generated by $N_0$ and the
corresponding subgroups for the other members of $B$. These selected
elements form an independent family in $S$.

Put $J=\gen{T,B}$. The group $J\cap N$ projects onto the first copy
of $S$. Conjugation by $N\leq S^m$ fixes each coordinate, so the transitive
coordinate action of $G$ is already that of $T_0$. Therefore
it projects onto every copy. Add the elements of $E\setminus B$
one by one. Each addition strictly enlarges the generated group, by
irredundance. Since the image modulo $N$ stays fixed, it strictly
enlarges its intersection with $N$. These intersections form a chain
of subdirect subgroups, with at most $m-1$ strict steps by
Lemma~\ref{lem:scott}. Thus
$|T|+|E|\leq\mu(G/N)+\mu'(S)+m-1$, as required.
\end{proof}

\section{Primitive block stabilizers}
\begin{lemma}\label{lem:primitiveindependent}
Let $L$ be primitive of degree $k\geq5$, with $L\ne A_k,S_k$.
Then $\mu'(L)\leq k-3$, except for the degree-six copy of $S_5$,
which has $\mu'=4$.
\end{lemma}
\begin{proof}
For $k\geq26$, the Mar\'oti and Whiston argument above bounds all
independent families, not only C-group markings, by $k-4$.
For $5\leq k\leq25$, the complete primitive catalogue and the
prime-factor bound leave exactly the following eleven entries:
\begin{equation}\label{eq:independentexceptions}
\begin{array}{c|l}
5& C_5:C_4\\
6& A_5,\ S_5\\
7& PSL_3(2)\\
8& AGL_3(2),\ PGL_2(7)\\
9& AGL_2(3),\ PSL_2(8):3\\
10& A_6:2^2\\
12& M_{12}\\
16& AGL_4(2).
\end{array}                                                    \end{equation}
The same required program checks this complete list and the exact
maximal-subgroup, normal-subgroup and quotient facts below.

The group $C_5:C_4$ has irredundant length at most two by its minimal
normal $C_5$ and cyclic prime-power quotient. Every proper subgroup
is cyclic or dihedral of order ten, hence also has bound two.
Whiston treats $A_5$ and $S_5$ in their natural degrees five.
Every maximal subgroup of $PSL_3(2)$ is $S_4$ or has at most two prime
factors, so a generating family has length at most four and a proper
subgroup has independent length at most three. The index-two extension
$PGL_2(7)$ consequently has independent length at most five.
Here and below we use the immediate subgroup version of
Lemma~\ref{lem:quotient}: $\mu'(G)\leq\mu'(G/N)+\mu'(N)$.

For $AGL_3(2)$ consider any subgroup $H$, its intersection $U$ with
the translation module $V$, and its linear image $P\leq GL_3(2)$.
If $\dim U=0$, the bound is four. If $\dim U=1$ or two, $P$ lies
in a line or plane stabilizer, each isomorphic to $S_4$ (its normal $2^2$ and quotient $GL_2(2)$
act as the full affine group on four points), so the bound
is $3+\dim U\leq5$. If $U=V$ and $P$ is irreducible, the normal
module costs one and the quotient costs at most four. If $P$ is
nontrivial and reducible, it lies in an $S_4$ stabilizer, and the
module costs at most two: three independent vectors would, by
Lemma~\ref{lem:simplequotients}, identify $V$ with a sum of three
one-dimensional, hence trivial, modules, contradicting the faithful
nontrivial action. Finally $P=1$ costs only three. Thus every $H$
has irredundant length at most five.

Every proper subgroup of $GL_2(3)$ lies in a maximal subgroup whose
order has at most four prime factors; the whole group has irredundant
length three by its Frattini quotient $S_4$. Thus $\mu'(GL_2(3))\leq4$,
and $\mu'(AGL_2(3))\leq2+4=6$. The earlier maximal-subgroup bound
for $PSL_2(8)$ gives independent length at most five, so its index-three
extension has bound six. The degree-ten group has normal $A_6$ and
quotient of order four, giving at most $4+2=6$. Every maximal subgroup
of $M_{12}$ has at most eight prime factors, so its independent length
is at most nine. Finally $AGL_4(2)$ has a four-dimensional normal
translation module and quotient $A_8$, giving $\mu'\leq4+6=10$.
All satisfy the asserted bounds.
\end{proof}

\section{The full imprimitive bound}
\begin{theorem}\label{thm:imprimitive}
Every transitive imprimitive permutation group of degree $n\geq9$
has maximum irredundant generating length at most $n-4$.
\end{theorem}
\begin{proof}
Choose a minimal nontrivial block, of size $k$, and let $m\geq2$
be the number of its translates, so $n=km$. The induced group $L$
of the setwise block stabilizer is primitive: a proper block for $L$
would propagate to a smaller block system for $G$. After choosing
identifications of the blocks by elements of $G$, the group embeds
in $L\wr S_m$. Its block image $P$ is transitive, and the projection
$K_1$ of the block kernel $K$ is normal in $L$; all other projections
are conjugate to it.

Blocks of size two, three or four are covered by
Theorems~\ref{thm:nine}, \ref{thm:twoblocks} and
Lemma~\ref{lem:smallbase}. Assume $k\geq5$.
If $m\geq3$ and $L<S_k$, then every subgroup of $L$ has independent
length at most $k-2$, by Whiston's equality clause. Successive
coordinate projections give $\mu'(K)\leq m(k-2)$, and therefore
\[
 \mu(G)\leq(m-1)+m(k-2)=n-m-1\leq n-4.
\]
If $m=2$ and $L\ne A_k,S_k$ is not the degree-six copy of $S_5$,
Lemma~\ref{lem:primitiveindependent} similarly gives
$\mu(G)\leq1+2(k-3)=n-5$.

It remains to treat $L=A_k,S_k$, and the exceptional degree-six
copy of $S_5$ with $m=2$. First take $L=A_k$ or $S_k$.
By simplicity of $A_k$ and its trivial centralizer in $S_k$, the normal subgroup $K_1$ is either trivial,
$A_k$, or $S_k$. In the trivial case $K=1$ and
$\mu(G)\leq m-1\leq n-4$. Otherwise
$N=K\cap A_k^m$ is subdirect in $A_k^m$: the commutator subgroup
of $K$ already projects onto $A_k$ in every coordinate. The quotient
$G/N$ embeds in $C_2\wr S_m$. Its maximum irredundant length is at
most $(m-1)+m=2m-1$ by the normal-module bound, so
Lemma~\ref{lem:simplekernel} and $\mu'(A_k)\leq k-2$ give
\[
 \mu(G)\leq (2m-1)+(k-2)+(m-1)=k+3m-4.
\]
For $k\geq5,m\geq3$ this is at most $km-4$, since
$(k-3)(m-1)\geq3$.
When $m=2$, the quotient lies in $C_2\wr S_2\cong D_8$;
every subgroup of $D_8$ has irredundant length at most two.
The sharper bound is then $\mu(G)\leq2+(k-2)+1=k+1\leq2k-4$.

Finally let $k=6,m=2,L\cong S_5$. Use the characteristic subgroup
$A_5\lhd L$ and the same argument in the abstract factors.
The block kernel is trivial or projects onto $A_5$ or $S_5$.
In the latter two cases Lemma~\ref{lem:simplekernel}, now with
$S=A_5$, gives $\mu(G)\leq2+3+1=6\leq12-4$.
This completes all cases and every parameter value.
\end{proof}

\begin{corollary}\label{cor:intransitive}
Let $(S_n;\rho_0,\ldots,\rho_{r-1})$ be a C-group with
$n\geq9$ and $r\in\{n-2,n-1\}$. Every maximal standard
parabolic is intransitive on the $n$ letters.
\end{corollary}
\begin{proof}
Each such parabolic has an independent distinguished family of
$r-1\geq n-3$ elements. The primitive bound and
Theorem~\ref{thm:imprimitive} exclude every transitive proper
subgroup except $A_n$. Suppose one parabolic is $A_n$ and let
$a$ be the omitted involution. For each other distinguished
involution $b$, the intersection property says
$A_n\cap\gen{a,b}=\gen{b}$. Since $aba\in A_n$, this forces
$aba=b$. Thus $a$ centralizes all of $A_n$ and
$S_n=A_n\times\gen{a}$, impossible. The centralizer of $A_n$
in $S_n$ is trivial for $n\geq4$, as is seen by commuting with
all three-cycles. Hence no maximal parabolic is transitive.
\end{proof}

